\section{Conclusion}

In Table~\ref{tab:model-comparison}, we summarize 
the final performance metrics for our two best models.

The \textbf{Regularized CNN model} achieved the lowest 
validation MSE and NLL, with a coverage of $66.8\%$, 
which is close to the expected $68\%$ for a well-calibrated 
Gaussian model. This indicates that the model not only 
provides accurate mass predictions but also delivers reliable 
uncertainty estimates.

In contrast, the \textbf{Trial and Error model} produced a 
very high coverage of $90.3\%$. While this might 
seem impressive, it can in be misleading: 
such high coverage indicates that the model's uncertainty 
estimates are very conservative, meaning it predicts error 
bars that are too wide. As a result, the model may appear to 
capture more true values within its predicted intervals, 
but this comes at the cost of less informative, less precise 
predictions. This interpretation is further supported by its 
higher validation MSE and NLL, which indicate that the Trial 
and Error model’s mass predictions are less accurate and its 
probabilistic predictions less well-calibrated.

Additionally, from the pull distributions, 
we can further assess the quality of the the models.
For the Regularized CNN model, 
the pull values are evenly distributed around zero, 
indicating well-calibrated uncertainties and no systematic
bias in the predictions. In contrast, the pull plot 
for the Trial and Error model exhibits an \textbf{S-shaped}
pattern, suggesting that its uncertainties are not well 
calibrated and may exhibit systematic deviations. This 
further supports our belief that, despite its higher
coverage, the Trial and Error model does not provide 
reliable uncertainty estimates.

Therefore, we conclude that the Regularized CNN model 
provides a better balance between prediction accuracy and 
uncertainty calibration for this task.

\begin{table}[ht]
    \centering
    \caption{Comparison of final performance metrics for the two best models.}
    \label{tab:model-comparison}
    \begin{tabular}{|l|c|c|c|}
        \hline
        \textbf{Model} & \textbf{Best Val NLL} & \textbf{Best Val MSE} & \textbf{Coverage (\%)} \\
        \hline
        Regularized CNN mode & $-0.08136$ & $0.073$ & $66.8$ \\
        Trial and Erro model & $0.0490$   & $0.217$ & $90.3$ \\
        \hline
    \end{tabular}
\end{table}

\textbf{Potential Improvements:}
The aforementioned work is by no means perfect and
can be improved in sevceral ways. First, repeating the 
hyperparameter optimization with Optuna using 
the negative log-likelihood (NLL) as the objective 
would directly target better uncertainty calibration.
Additionally, 
experimenting with more advanced architectures—such as 
hybrid models that combine Convolutional Neural Networks 
(CNNs) with Transformers, or using Vision Transformers 
(ViTs) and attention-based mechanisms—could further enhance
performance on image data. 

Another potential direction involves a more detailed investigation 
of the auxiliary input features. In this study, we included the full 
range of redshift and plasma temperature ($kT$) values as input features. 
Future work could analyze the distributions of these features in more 
detail, potentially identifying specific ranges where the model's 
performance is stronger or where the data is more reliable. 


